\DocumentMetadata{
  pdfversion=2.0,
  pdfstandard=ua-2,
  lang=en-US,
  tagging=on,
  tagging-setup={math/setup=mathml-SE,tabsorder=structure}
}
\documentclass[11pt,letterpaper]{article}
\usepackage[margin=0.68in,includefoot]{geometry}
\usepackage{newcomputermodern}
\usepackage{amsmath,amscd,mathtools}
\usepackage{tikz,adjustbox}
\providecommand{\square}{\mdlgwhtsquare}
\usepackage[hidelinks]{hyperref}
\hypersetup{
  pdftitle={Analysis First-Year Exam, May 2017, Part 2},
  pdfauthor={University of Florida Department of Mathematics},
  pdfsubject={Graduate mathematics examination},
  pdfdisplaydoctitle=true,
  bookmarksopen=true
}
\setcounter{secnumdepth}{0}
\setlength{\parindent}{0pt}
\setlength{\parskip}{0.28em}
\setlength{\emergencystretch}{3em}
\setlength{\leftmargini}{2.15em}
\setlength{\leftmarginii}{2.65em}
\setlength{\leftmarginiii}{3em}
\renewcommand{\labelenumi}{\textbf{\theenumi.}}
\pagestyle{empty}
\raggedbottom
\allowdisplaybreaks
\newenvironment{examproblems}[1]{%
  \begin{enumerate}%
  \setcounter{enumi}{#1}%
  \setlength{\topsep}{0.35em}%
  \setlength{\partopsep}{0pt}%
  \setlength{\itemsep}{0.48em}%
  \setlength{\parsep}{0.18em}%
}{\end{enumerate}}
\newenvironment{examsubparts}{%
  \begin{enumerate}%
  \setlength{\topsep}{0.18em}%
  \setlength{\partopsep}{0pt}%
  \setlength{\itemsep}{0.16em}%
  \setlength{\parsep}{0.1em}%
}{\end{enumerate}}
\newenvironment{examinstructions}{%
  \par\begingroup\itshape
}{%
  \par\endgroup\vspace{0.18em}
}
\makeatletter
\renewcommand\section{\@startsection{section}{1}{0pt}%
  {-1.25ex plus -.25ex minus -.1ex}%
  {0.55ex plus .1ex}%
  {\normalfont\large\bfseries}}
\renewcommand{\@maketitle}{%
  \newpage\null\vskip -1em%
  {\footnotesize\bfseries
    \href{https://www.ufl.edu/}{University of Florida}, \href{https://math.ufl.edu/}{Department of Mathematics}\par}%
  \vskip -0.5em%
  \tagstructbegin{tag=Title}%
    {\LARGE\bfseries\href{https://gma.math.ufl.edu/exams/analysis-first-year/}{Analysis First-Year Exam}, May 2017, Part 2\par}%
  \tagstructend%
  \vskip 0.25em%
  \tagmcbegin{artifact}%
    \hrule height 1pt%
  \tagmcend%
  \vskip 1em%
}
\makeatother
\title{Analysis First-Year Exam, May 2017, Part 2}
\author{}
\date{}
\begin{document}
\maketitle
\thispagestyle{empty}
\begin{examinstructions}
Answer FOUR questions in detail. State carefully any results used without proof.
\end{examinstructions}
\begin{examproblems}{0}
\item
Let \((f_n)_{n=0}^{\infty}\) be a sequence of differentiable functions. Decide whether each implication is valid, giving proof or counterexample as appropriate:
\begin{examsubparts}
\item[\textnormal{(i)}]
if \(f_n \to f\) uniformly on \((−\infty,\infty)\) then \(f_n^2 \to f^2\) uniformly on \((−\infty,\infty)\);
\item[\textnormal{(ii)}]
if \(f_n \to f\) uniformly on \([−1,1]\) then \(\int_{−1}^1 f_n \to \int_{−1}^1 f\);
\item[\textnormal{(iii)}]
if \(f_n \to f\) uniformly on \((−1,1)\) then~\(f\) is differentiable and \(f_n' \to f'\) uniformly on \((−1,1)\).
\end{examsubparts}
\item
Let~\(f\) be a continuous real-valued function on \([0,1]\) and let \(\alpha \in (0,\infty)\). Assume that 
\[
\int_0^1 t^{n\alpha}f(t)\,dt=0
\]
 for all but finitely many values of \(n\in\mathbb N\). What conclusions can be drawn about~\(f\)?
\item
Let \(\mathcal F\) be an equicontinuous family of real-valued functions on the compact metric space~\(X\). Denote by \(A\subseteq X\) the set whose elements are precisely those \(a\in X\) at which \(\mathcal F\) is bounded in the sense that \(\{f(a):f\in\mathcal F\}\subseteq\mathbb R\) is bounded. Prove that~\(A\) is both open and closed.
\item
Let \((f_n)_{n=0}^{\infty}\) be a sequence of measurable real-valued functions. Decide whether each of the following sets is measurable:
\begin{examsubparts}
\item[\textnormal{(i)}]
\(\{\omega:(f_n(\omega))_{n=0}^{\infty}\text{ is unbounded}\}\);
\item[\textnormal{(ii)}]
\(\{\omega:(f_n(\omega))_{n=0}^{\infty}\text{ is periodic}\}\);
\item[\textnormal{(iii)}]
\(\{\omega:(f_n(\omega))_{n=0}^{\infty}\text{ has distinct terms}\}\).
\end{examsubparts}
\item
Let \((\Omega,\mathcal A,\mu)\) be a measure space that is finite in the sense that \(\mu(\Omega)<\infty\). Let \((f_n)_{n=0}^{\infty}\) be a sequence of non-negative measurable functions converging pointwise to~\(f\) on~\(\Omega\). True or false (proof or counterexample): 
\[
\int_\Omega \frac{1}{1+f_n}\,d\mu \to \int_\Omega \frac{1}{1+f}\,d\mu \quad\text{as }n\to\infty.
\]
\end{examproblems}
\end{document}
